\scanpage{46}
\mspage{44}

ou \unsure{du moins|disons} (équivalences faibles) en bijections, \uline{ssi} l'hom précédent est un isom. de ind-objets de Ind Cat hot.

Ce sera donc cela la condition naturelle à imposer aux \deleted{\unsure{catég}} ind-catégories de chemins

\[
 \underline{\underline{\mathrm{Ch}}}_{\Sigma}(M,\unsure{a,b|a;b})
 \qquad(a,b\in\mathrm{Ob}\,M)
\]

provenant d'une ``catégorie de modèles'' $M,\Sigma$.

Bien sûr, disposant d'un système inductif explicite de catégories (pas seulement d'objets de Cat hot)

\[
 K_n=\underline{\mathrm{Ch}}_{\Sigma,n}(M;a,b),
\]

il est tentant de passer à la limite inductive

\[
 K_\infty=\varinjlim_n\underline{\mathrm{Ch}}_{\Sigma,\deleted{\unsure{n|\infty}}}(M;a,b)
 \overset{\mathrm{def}}{=}\underline{\mathrm{Ch}}_\Sigma(M;a,b),
\]

et ceci d'autant plus que les foncteurs de transition sont des monomorphismes et pleinement fidèles, de sorte que \deleted{$K_{\unsure{\infty|n}}$ est} les $K_n$ s'identifient à une suite de souscatégories pleines de $K_\infty$, dont la réunion est $K_\infty$. Mais d'autre part je ne sais donner aucune interprétation \unsure{convaincante} de $K_\infty$, \unsure{comme|connexe} \deleted{de} de celle du ind-objet \inserted{de Cat hot} défini par $\underline{\underline K}=(K_n)_{n\in\mathbb N}$, qui ind-représente le foncteur remarquable

\[
 H=H_{\unsure{M,\Sigma|M},a,b}:\mathrm{Cat\,hot}^{\unsure{o|0}}
 \unsure{\relbar\joinrel\relbar|\longrightarrow}\mathrm{Ens}
\]

factorisant le foncteur

% archive sheet number: 45
